\section{从语言模型到原生多模态：问题到底变了什么}

这堂课不是简单罗列“能看图、能说话”的产品，而是追问一个更基础的系统问题：\textbf{LLM 的成功究竟有多少来自可迁移的 sequence-modeling 原理，又有多少依赖语言这种特殊的信息载体？} Victoria Lin 先把 language model 的成功拆成 next-token prediction、规模化数据、instruction following、reasoning 与 planning，再逐步加入 image、audio、video 和 physical interaction。读者需要始终区分三层目标：接受多模态输入、生成非文本输出、在物理世界中持续感知与行动。三者不是同义词，也不应由同一张产品截图推断。

\begin{warningbox}{来源边界：本讲不代表 Thinking Machines 的内部实现}
讲者在 00:00:59--00:01:12 明确声明，课堂只讨论 public materials，观点也不代表雇主。本讲义因此只使用 Stanford 官方 deck、公开录像、manual captions 与课堂日前公开论文；不会把 deck 中的概念外推成 Thinking Machines 的未公开系统细节。
\end{warningbox}

\subsection{先复盘 LLM scaling baseline}

在进入多模态之前，先看语言模型为什么成为一个可迁移的设计模板。基础接口很简单：把离散文本转成 token 序列，使用 Transformer 对条件分布建模，再通过预训练、post-training 和 inference-time computation 获得知识、指令遵循、工具使用与规划能力。关键不是每项能力都由 next-token prediction 自动保证，而是这个统一接口让数据、模型、优化与部署可以共同规模化。

\lecturefigure{slide-005.jpg}{语言模型从问答、代码到 agent 的能力栈}{官方 deck 第 5 页；课堂 00:01:20--00:02:52}

\begin{knowledgebox}{读图：不要把应用图标当成机制证明}
这张图把 ChatGPT、coding assistant 与 agent-style workflow 放在同一页，是为了说明语言接口已从“回答问题”扩展到“执行多步任务”。先比较中间的共同核心：所有应用都把复杂目标压成符号序列，再让模型预测、调用工具或迭代修正。图中没有证明 emergent capability 必然来自某个单一组件；它提供的是后续迁移问题的 baseline：多模态系统能否保留同样统一、可扩展的接口？
\end{knowledgebox}

\teachervoice{00:01:20--00:02:52，老师把知识、instruction following、chain-of-thought 与 planning 描述为规模化语言模型逐步显现的能力。讲义提醒：这里是经验性能力栈，不是“只要增加参数就必然出现所有能力”的定理。}

\subsection{为什么 language modeling alone 不够}

语言高度压缩了人类经验，但互联网和物理世界并不只由文字构成。图像携带空间结构，音频携带时间、说话人和环境信号，视频携带变化与因果线索，机器人还要把感知转成连续控制。若系统只能把这些信息先交给外部工具，再让 LLM 阅读一段文字摘要，它仍可能是强大的 orchestration layer，却还不是本讲所谓的 native multimodal model。

\lecturefigure{slide-006.jpg}{从“语言足够”转向“语言不够”的核心命题}{官方 deck 第 6 页；课堂 00:02:56--00:03:00}

\begin{importantbox}{本讲的 native multimodality 工作定义}
原生多模态系统把不同 modality 的表示、序列位置、训练目标或参数路径纳入同一模型设计，使模型能够联合感知、条件化、推理，并在需要时直接生成非文本模态。它不是“给文本模型接一个 OCR/API”这么宽泛，也不要求所有 modality 必须共享完全相同的 encoder、loss 或参数。
\end{importantbox}

这一页只有一句话，却承担了全讲最重要的转折。它不否定 language model，而是把问题从“怎样扩大语言模型能力”改写为“怎样让 sequence model 处理语言之外的高带宽信号，同时保留规模化训练的优点”。

\lecturefigure{slide-007.jpg}{数字世界的信息处理与物理世界的多模态经验}{官方 deck 第 7 页；课堂 00:03:00--00:04:18}

\begin{knowledgebox}{读图：左半边与右半边的难度不对称}
左侧 digital world 中的文档、图表、代码、网页与音频通常已被切成可离线处理的对象；模型可以等待完整输入，再生成回答。右侧 physical world 则要求持续感知、低延迟反应、空间一致性和动作闭环。图中最值得比较的不是“有多少 modality”，而是 information arrival、latency budget 与 consequence：漏读一张图表通常可重试，机器人在真实环境中的错误动作可能不可逆。
\end{knowledgebox}

\teachervoice{00:03:23--00:03:49，老师强调若 AI 要与人实时互动，就必须直接处理 image、audio、video 等信号。这里的“实时”不仅是 streaming UI，而是感知、状态更新和响应都处在环境时间约束内。}

\subsection{现有系统位于哪一层}

市场上常把“multimodal”作为统一标签，但工程上至少应问四个问题：输入是否多模态？输出是否多模态？表示是否联合训练？模型是否在实时物理闭环中运行？同一个系统可能在第一项很强、第二项有限、后两项仍依赖外部组件。

\lecturefigure{slide-008.jpg}{课堂列举的 native multimodal model 语境}{官方 deck 第 8 页}

\begin{warningbox}{读图时不要进行产品能力倒推}
Gemini、Kimi、Qwen 等名称用于说明 2026 年已有多模态产品与研究系统；这张图不是统一 benchmark，也没有证明它们采用相同 architecture、loss 或 modality coverage。后文的技术结论应来自机制图和论文，而不是根据品牌名称猜测内部实现。
\end{warningbox}

\subsection{本章小结}

LLM 提供了一个可规模化的 sequence-modeling baseline，但 native multimodality 增加了三类压力：representation 不再天然离散，output objective 不再只有 categorical next-token loss，physical interaction 还引入实时性与行动后果。后续所有模型都可以放进同一个设计空间：\textbf{如何表示、如何排序、预测什么、哪些参数共享、哪些能力真正正迁移。}

